\section{低维结果的证明}
\label{low:appendix}

本附录给出第~\ref{low:section}节的刚性与弯曲路径证明，并用不变量区分指定四比特基准测量与已知大半径参数段。

\subsection{\texorpdfstring{$n=3$}{n=3} 时刚性的精确证书}
\label{low:app-rigidity}

记 $b_z=i^{f_0(z)}$，并取 $\theta_0=0$。对 带符号 物理 Pauli 算符
$\widetilde P$，定义固定矩阵系数
\begin{equation}
 K_P(z,w)=\overline{b_w}b_z
 \langle w|C_n^\dagger\widetilde P C_n|z\rangle.
 \label{low:finite-coeff}
\end{equation}
种子分量及其各阶导数均由下列有限恒等式给出：
\begin{equation}
 2^n m_P(\theta)=\sum_{z,w}K_P(z,w)e^{i(\theta_z-\theta_w)}.
 \label{low:finite-sum}
\end{equation}
特别地，
\begin{equation}
 2^n\partial_lm_P(0)
 =i\sum_{z,w}K_P(z,w)
 (\mathbf1_{z=l}-\mathbf1_{w=l}).
 \label{low:jacobian-formula}
\end{equation}
上述公式直接采用原电路，明确保留其符号与归一化。

对 $n=3$，令 $\mathcal F$ 由八个差 $m_P-m_{X_1}$ 组成，行顺序为
\begin{equation*}
 P=Y_1,Z_1,X_2,Y_2,Z_2,X_3,-Y_3,Z_3.
\end{equation*}
列对应 $\theta_1,\ldots,\theta_7$。由式~\eqref{low:jacobian-formula} 得
\begin{equation}
 4D\mathcal F(0)=
 \begin{pmatrix}
 -2&-2&0&1&1&0&2\\
 -1&-2&0&-1&2&2&0\\
 -1&-1&-1&0&1&0&1\\
 -1&-1&-1&0&1&1&0\\
 -1&-2&0&1&0&0&2\\
 -2&-1&0&1&0&0&2\\
 0&-2&-1&1&2&0&0\\
 -1&0&-2&1&0&2&0
 \end{pmatrix}.
 \label{low:jacobian-three}
\end{equation}
所示整数矩阵前七行的行列式为 $4$，故对应导数的行列式为
$4/4^7=1/4096$。对这七个实解析分量差应用逆函数定理，得到其共同零点的局部唯一性。
指定点附近的每个正四面体测量都满足这些方程：其全部 带符号 分量保持正号，
故分量平方相等蕴含分量本身相等。这就在完整七维相位空间中证明了命题~\ref{low:rigidity}。

\subsection{\texorpdfstring{$n=4$}{n=4} 时的全部实相位直线}
\label{low:app-lines}

现在令 $\mathcal F$ 由其余十一个 带符号 分量与 $X_1$ 的差组成，算符顺序为
\begin{equation*}
 Y_1,Z_1,X_2,Y_2,Z_2,X_3,Y_3,Z_3,X_4,-Y_4,Z_4.
\end{equation*}
对直线 $\theta=td$，由式~\eqref{low:finite-sum} 得到精确导数公式
\begin{equation}
 16\frac{d^k}{dt^k}m_P(td)\bigg|_{t=0}
 =\sum_{z,w}K_P(z,w)\bigl[i(d_z-d_w)\bigr]^k.
 \label{low:line-derivative}
\end{equation}
一阶方程的完整解由下列四个任意实参数 $a,b_*,c,\kappa$ 给出：
\begin{equation}
 \begin{aligned}
 d_0&=0,\qquad d_1=d_2=d_9=a,\\
 d_3=d_8&=a-c,\qquad d_{10}=b_*-c,\\
 d_4=d_5&=-a+b_*-c/2+\kappa,\\
 d_6=d_7=d_{12}=d_{13}&=c,\\
 d_{11}&=b_*,\qquad d_{14}=d_{15}=\kappa.
 \end{aligned}
 \label{low:kernel-four}
\end{equation}
这里给出式~\eqref{low:jacobian-formula} 的精确线性代数核验：
所示四个方向均被 $D\mathcal F(0)$ 消去，且 $16D\mathcal F(0)$
的列 $(1,2,3,4,5,6,7,8,10,12,14)$ 组成的子式行列式为 $65536$。
因此其秩为 $11$，式~\eqref{low:kernel-four} 穷尽了线性化核。

令 $G_P(t)=16[m_P(td)-m_{X_1}(td)]$。
由式~\eqref{low:line-derivative} 得到的二阶条件包括
\begin{align}
 G''_{Y_1}(0)&=-2c(2a-c),\nonumber\\
 G''_{Y_2}(0)&=4c(a-2b_*+2c).
 \label{low:second-order}
\end{align}
若 $c\ne0$，两式为零迫使 $a=c/2$、$b_*=5c/4$。
下一个条件变为
\begin{equation}
 G''_{X_2}(0)=-\frac{c(c+4\kappa)}2=0,
\end{equation}
从而 $\kappa=-c/4$。但此时相同的导数公式给出
\begin{equation}
 G'''_{Y_1}(0)=-\frac{39c^3}{8}\ne0.
 \label{low:third-order}
\end{equation}
因此必有 $c=0$。剩余必要二阶条件为
\begin{equation}
 G''_{X_2}(0)=-8(a-b_*)(a-b_*-\kappa)=0.
 \label{low:two-planes}
\end{equation}
定义三个整数输入标号集合
\begin{align}
 \mathcal A&=\{1,2,3,8,9,10,11\},\nonumber\\
 \mathcal B&=\{1,2,3,8,9,14,15\},\nonumber\\
 \mathcal C&=\{4,5,14,15\}.
 \label{low:line-supports}
\end{align}
分支 $b_*=a$ 给出 $d=a\mathbf1_{\mathcal A}
+\kappa\mathbf1_{\mathcal C}$；分支 $\kappa=a-b_*$ 给出
$d=b_*\mathbf1_{\mathcal A}+(a-b_*)\mathbf1_{\mathcal B}$。
将两分支代入有限和~\eqref{low:finite-sum}，对每个 带符号 局部 Pauli 算符均得到
\begin{equation}
 16m_P(td)=e^{iat}+e^{-iat}.
 \label{low:planes-sufficient}
\end{equation}
这证明了全部实参数上的充分性，而不止有限阶的一致。
因此全部可行方向恰为
\begin{equation}
 \operatorname{span}_{\mathbb R}
 \{\mathbf1_{\mathcal A},\mathbf1_{\mathcal C}\}
 \ \cup\ 
 \operatorname{span}_{\mathbb R}
 \{\mathbf1_{\mathcal A},\mathbf1_{\mathcal B}\}.
 \label{low:line-classification}
\end{equation}
基点处的正号使局部规则性蕴含上述 带符号 方程。
任意实方向上的分量差都是有限指数和，故在开区间上成立的等式延拓至整条直线。
式~\eqref{low:planes-sufficient} 与~\eqref{low:line-classification}
证明了定理~\ref{low:theorem} 的第一部分，其中没有限制 $d$ 所取值的个数。

\subsection{弯曲分支及其完整单调区间}
\label{low:app-curve}

式~\eqref{low:masks} 指定式~\eqref{low:curve} 的全部相位。
输入以最高位在左的整数 $z=0,\ldots,15$ 编号。
将 $\theta_z=bu_z+\delta v_z$ 代入式~\eqref{low:finite-sum}，
恰得到式~\eqref{low:two-components} 中的两种表达式。
这些恒等式对独立实参数 $b,\delta$ 成立，无须预先选取根分支。

\begin{equation}
 \begin{array}{c|rrrrrrrr}
 z&0&1&2&3&4&5&6&7\\
 u_z&0&1&1&1&2&2&0&0\\
 v_z&0&1&1&0&1&1&1&1\\[2pt]
 z&8&9&10&11&12&13&14&15\\
 u_z&1&1&0&0&0&0&1&1\\
 v_z&0&1&0&1&1&1&0&0
 \end{array}
 \label{low:masks}
\end{equation}

令 $q=-\delta$、$z=\tan(q/2)$；仅在以下根分支计算中，字母 $z$ 表示实变量。
半角公式给出
\begin{equation}
 H-T=\frac{2(Az^2-Bz+C)}{1+z^2}.
 \label{low:quadratic-identity}
\end{equation}
下面验证式~\eqref{low:delta} 的根及其完整定义域。
令 $x=\tan(b/2)$、$k=\sqrt2-1$，则
\begin{equation}
 Az^2-Bz+C=\frac{x+1}{(1+x^2)^2}F(x,z),
 \label{low:halfangle-polynomial}
\end{equation}
其中
\begin{align}
 F(x,z)={}&-x^3z^2+x^3z-3x^2z^2-3x^2z-4x^2\nonumber\\
 &+3xz^2+xz+4x+z^2-3z.
 \label{low:polynomial}
\end{align}
在正方形 $0\le x,z\le k$ 上，
\begin{equation}
 F_{xx}=6z[x(1-z)-(1+z)]-8<0.
\end{equation}
利用 $k^2+2k-1=0$ 及 $4-12k<0$，得到
\begin{align}
 F_x(x,z)&\ge F_x(k,z)
 =4-8k+(4-12k)z\nonumber\\
 &\ge20k-8=20\sqrt2-28>0.
 \label{low:fx-positive}
\end{align}
另一偏导数满足
\begin{align}
 F_z&=2z(1+3x-3x^2-x^3)+x^3-3x^2+x-3\nonumber\\
 &\le2k(1+3k)+k^3+k-3
 =5-4\sqrt2<0.
 \label{low:fz-negative}
\end{align}
边界值满足
\begin{align}
 F(x,x)&=-x(1+x^2)(x^2+2x-1)\ge0,\nonumber\\
 F(x,k)&\le F(k,k)=0.
 \label{low:boundary}
\end{align}
因此每个 $x\in[0,k]$ 都有唯一根 $z(x)\in[x,k]$；
$z(0)=0$、$z(k)=k$，且在区间内部有严格不等式 $x<z(x)<k$。
根处满足 $F_z<0$。由于 $A>0$，对应二次式的判别式严格为正，且此根为较小根。
有理化根公式得到
\begin{equation}
 z(x)=\frac{2C}{B+\sqrt{B^2-4AC}},
 \label{low:root-identified}
\end{equation}
该式在 $x=0$ 处也有良好定义。由此验证了所给闭式公式，而非仅以隐式存在性来定义路径。

对已经识别的单根求导，得到 $z'(x)=-F_x/F_z>0$，故
$q'(b)>0$ 且 $b\le q(b)\le\pi/4$。
将 $T$ 写为 $T=\cos(b-q)+\sin q$，则
\begin{equation}
 T'(b)=\sin(q-b)+q'(b)\,[\cos q-\sin(q-b)]>0.
 \label{low:monotone}
\end{equation}
方括号严格为正：若取等号，必须同时有 $q=\pi/4$ 和 $b=0$，与已确定的端点值矛盾。
第一项非负。最后，由 $T(0)=1$ 及 $T(\pi/4)=1+1/\sqrt2$
得到式~\eqref{low:curve-range} 的精确像。在原点处，
$z'(0)=4/3$、$q'(0)=4/3$、$r'_4(0)=\sqrt3/6$。

在式~\eqref{core:bloch-orbit} 中取 $n=4$、$m=T/8$，即可得到全部约化态，
每个固定正四面体方向各出现四次。式~\eqref{low:curve} 的酉分解给出正交完备测量，
且 $b=0$ 恢复指定的 $D_{f_0}$。

\subsection{基准测量与已知 \texorpdfstring{$n=4$}{n=4} 参数段的非等价性}
\label{low:app-invariant}

本节证明，指定基准测量与文献~\cite{PauwelsGisin2026} 附录 C 所给 $n=4$ 参数段中的任何测量，均不能通过局域酉变换、比特置换及结果重标号相互转换。
对 $a\in[0,\pi/3]$，令 $x=e^{ia}$，将其输入相位写为 $i^{f_z}x^{k_z}$。
按 $z=0,\ldots,15$ 排序，两个固定向量为
\begin{align}
 f={}&(0,0,0,3,0,0,1,0,\nonumber\\
 &\hspace{12mm}0,0,0,2,0,1,3,3),\nonumber\\
 k={}&(0,1,-1,-1,0,1,0,0,\nonumber\\
 &\hspace{12mm}-1,2,0,1,3,3,1,0).
 \label{low:f159-data}
\end{align}
这里 $a=\pi(1-\tau)/3$，$\tau$ 为该文使用的参数。
种子 $C_4(i^{f_z}x^{k_z})_z/4$ 的单比特本征值为
$(1\pm\sqrt3\,c(x))/2$，其中
\begin{equation}
 c(x)=\frac{x^3+x^{-3}+2x+2x^{-1}}{16}
 =\frac{\cos3a+2\cos a}{8}.
 \label{low:f159-component}
\end{equation}
该式在所给参数段上非负，从 $3/8$ 递减至零，给出已知端点 $3r_E(4)$。

对$n=4$ 的纯态定义
\begin{equation}
 S(\psi)=\operatorname{Tr}(\rho_{12}^{2})
 +\operatorname{Tr}(\rho_{13}^{2})
 +\operatorname{Tr}(\rho_{14}^{2}).
 \label{low:purity-invariant}
\end{equation}
该和在局域酉及比特置换下不变：纯态的互补子系统具有相同非零谱，
故置换只会重新排列这三个无序 $2|2$ 划分项。
Pauli 轨道测量的每一列均与种子共享这些不变量，因此结果重标号不能改变比较结果。
在指定 $f_0$ 点，三个纯度依次为 $(63,39,39)/128$，故 $S_0=141/128$。

等价测量必须具有相同的单比特谱和 $S$。
在此参数段上，两条件分别要求 $c=1/8$ 与 $S=141/128$。
由式~\eqref{low:f159-data} 的振幅计算各约化态的纯度，并利用 $x\ne0$ 清除分母，
得到必要方程 $P(x)=Q(x)=0$，其中
\begin{align}
 P(x)={}&x^6+2x^4-2x^3+2x^2+1,\nonumber\\
 Q(x)={}&6x^{16}-3x^{14}+74x^{12}+48x^{10}\nonumber\\
 &+134x^8+48x^6+74x^4-3x^2+6.
 \label{low:invariant-polynomials}
\end{align}
Euclid 算法排除了共同根，其第一步余式为
\begin{equation}
 \begin{aligned}
 Q\bmod P={}&816x^5-2340x^4+2168x^3\\
 &-1240x^2+616x-596,
 \end{aligned}
\end{equation}
其与 $P$ 的最大公因式为 $1$。因此，原参数段中没有任何测量在所允许的操作下等价于指定的 $f_0$ 测量。




